% =====================================================================
% PAPER II ("The Strike")
% Exceptional eigenvalues of Bianchi orbifolds by verified finite
% elements. II: Congruence levels of the Picard group
% ---------------------------------------------------------------------
% OUTLINE (compilable skeleton).  Content notes in [[double brackets]]
% map to repository sources: CONGRUENCE.md, 0cuspchecks.md,
% congruence_prototype.py, m3p_certify.py (worktree np9 versions),
% theorem4postmortem.md.
% Reviewer burden by design: the computational engine is Paper I's and
% is cited, not re-derived.  Reviewers here evaluate (i) the
% reference-cell principle, (ii) the exact coset gluing, (iii) the new
% multi-cusp geometry.  Nothing numerical-analytic is new except the
% per-cusp bookkeeping of Theorem G1p.
% =====================================================================
\documentclass[11pt]{amsart}
\usepackage{amsmath,amssymb,amsthm,booktabs}
\newtheorem{theorem}{Theorem}[section]
\newtheorem{lemma}[theorem]{Lemma}
\newtheorem{proposition}[theorem]{Proposition}
\newtheorem{remark}[theorem]{Remark}
\newcommand{\HH}{\mathbb{H}^3}
\newcommand{\ZZ}{\mathbb{Z}}
\newcommand{\p}{\mathfrak{p}}
\newcommand{\note}[1]{\par\smallskip\noindent\textbf{[[#1]]}\par\smallskip}

\title[Verified exclusion of exceptional eigenvalues, II]{Exceptional
eigenvalues of Bianchi orbifolds by verified finite elements.\\
II: Congruence levels of the Picard group}
\author{}
\date{Draft outline, 2026-07-11}

\begin{document}
\begin{abstract}
Using the verified finite-element exclusion machinery of Paper I, we
prove that the Hecke congruence quotients
$\Gamma_0(\p)\backslash\HH^3$ of the Picard group, for the Gaussian
primes $\p$ of norm $5$, $9$ and $13$, have no Laplace eigenvalue in
$(0,1)$.  These quotients lie beyond the reach of trace-formula
positivity, whose margin degrades linearly with volume.  The key
structural input is a reference-cell principle: the congruence quotient
is tiled by exactly $N\p+1$ isometric copies of the level-one cell,
glued by exact arithmetic over $\ZZ[i]/\p$, with both cusp collars
assembling as products at the same truncation height --- so every
interval-arithmetic error bound is inherited from level one unchanged,
and only the size of the final verified linear algebra grows with the
index.
\note{Abstract: lead with the theorems and the ``unreachable by B<1''
point; then the principle; end with reproducibility.}
\end{abstract}
\maketitle

% ---------------------------------------------------------------------
\section{Introduction}
\note{Sources: CONGRUENCE.md header + \S8; parent README roadmap.}
\begin{itemize}
\item The Selberg-type problem at congruence level over $\mathbb{Q}(i)$;
  why level one is validation and $\Gamma_0(\p)$ is open territory.
\item The volume obstruction to trace-formula positivity ($B\sim$
  vol-linear; $B\approx 1.9$ already at $N\p=5$) vs.\ this method's
  margin, which tracks the spectral gap: the measured ladder
  $\mu \approx 7.3,\ 4.4,\ 2.4,\ 1.45$ for $N\p = 1, 5, 9, 13$
  (sublinear).  [[Table.]]
\item Main theorems (below); method summary in one paragraph with all
  engine citations to Paper I.
\item What is genuinely new here: Facts A--B (exact tiling and collar
  assembly), the conforming/relaxed gluing dichotomy, the multi-cusp
  criterion, Theorem G1p, and the two Rump refinements that the
  $N\p=13$ certificate forced (stated in Paper I, stress-tested here).
\end{itemize}

\begin{theorem}[$N\p=5$]\label{thm:5}
$\Gamma_0((2+i))\backslash\HH^3$ has no Laplace eigenvalue in $(0,1)$.
\end{theorem}
\begin{theorem}[$N\p=9$]\label{thm:9}
$\Gamma_0((3))\backslash\HH^3$ has no Laplace eigenvalue in $(0,1)$.
\end{theorem}
\begin{theorem}[$N\p=13$]\label{thm:13}
$\Gamma_0((3+2i))\backslash\HH^3$ has no Laplace eigenvalue in $(0,1)$.
\end{theorem}

% ---------------------------------------------------------------------
\section{The multi-cusp criterion}\label{sec:crit}
\note{Sources: CONGRUENCE.md \S2; 0cuspchecks.md C1, C4, C5, C6.}
\begin{itemize}
\item Cusps of $\Gamma_0(\p)$ (prime level: $\infty$ and $0$); scaling
  matrices $\sigma_\alpha\in\Gamma$; conjugated stabilizers (C1:
  $\sigma_0^{-1}(a,0;c,d)\sigma_0=(d,-c;0,a)$, lattice $\Lambda_0=\p$,
  fold present).
\item Per-cusp data table: $\Lambda_\alpha$, $|T_\alpha|$
  ($=1/2,\ N\p/2$), shortest dual vector ($1,\ 1/\sqrt{N\p}$), mode
  bounds $4\pi^2Y^2/N\p$ (valid for $N\p<4\pi^2Y^2\approx 61.7$).
\item Zero-mode defects and the constraint: $\beta_\infty$,
  $\beta_0=2(1-s)/(N\p\,Y^2)$; equal truncation heights $\Rightarrow$
  shared $\kappa_c$ and a rank-one constraint in $t_\infty+t_0$.
\item Precise invariance of both collars inherited from level one via
  $\sigma_\alpha\in\Gamma$ (C6) — nothing about $\Gamma_0(\p)$ used
  beyond inclusion.
\item Criterion theorem (multi-cusp) and pencil form; proof by the
  per-cusp copies of Paper I's Lemmas 1--3.
\end{itemize}

% ---------------------------------------------------------------------
\section{The reference-cell principle}\label{sec:refcell}
\note{Sources: CONGRUENCE.md \S1; 0cuspchecks.md C2, C3.  This is the
conceptual heart of the paper; give it room.}
\begin{itemize}
\item \textbf{Fact A} (exact tiling): right cosets
  $\Gamma_0(\p)\backslash\Gamma\cong\mathbb{P}^1(\ZZ[i]/\p)$;
  $D=\bigcup\gamma_i F$ with all copies isometric to the reference
  cell; per-copy element data equals reference data in pulled-back
  coordinates.
\item \textbf{Fact B} (collar assembly): representatives
  $\gamma_k=(0,-1;1,k)$; the computation
  $\sigma_0^{-1}\gamma_k=(1,k;0,1)$ is height-preserving, so the
  cusp-$0$ chimneys assemble to $T_0\times(Y,\infty)$ at the
  \emph{same} height $Y$.  Fold-compatibility of the tiling
  $T_0=\bigsqcup_k(T+k)$: the coset identity $k\equiv\varepsilon\lambda
  \pmod\p$ (C2) — the one step that is not automatic in the semidirect
  product.
\item \textbf{The Principle}: every interval-enclosed quantity is a
  level-one per-cell quantity or exact integer combinatorics; error
  radii are independent of the index; only the verified linear algebra
  grows.  [[State as a proposition about the certificate, with the
  volume identity $(N\p+1)\,\mathrm{vol}(K)+\sum|T_\alpha|/(2Y^2)=
  (N\p+1)\,\mathrm{vol}(F)$ as the sanity check.]]
\end{itemize}

% ---------------------------------------------------------------------
\section{Exact gluing and the symmetric mesh}\label{sec:glue}
\note{Sources: CONGRUENCE.md \S\S3--4; congruence\_prototype.py
(GaussianResidueRing); 0cuspchecks.md C7.}
\begin{itemize}
\item Side pairings of the reference cell and their face maps; the
  right action on $\mathbb{P}^1(\ZZ[i]/\p)$; gluing rule
  $j=(c{:}d)\cdot\delta^{-1}$ with derivation.
\item The dichotomy: self-identifications are relaxed (valid, enlarges
  the test space); interfaces between distinct copies are interior and
  must be conforming — relaxing them would hand each copy its own
  constant (the criterion would correctly fail).
\item The 24-split mesh: invariance under all face maps (mirror maps
  rule out Kuhn diagonals); Lemma G on quarter-triangles.
\item Residue arithmetic: $\mathbb{F}_5$ ($i\mapsto 3$),
  $\mathbb{F}_9$ (pairs mod $3$ — not a prime field),
  $\mathbb{F}_{13}$ ($i\mapsto 8$); the assert suite (permutations,
  round-trip, class sizes $\{1,N\p\}$, connectivity,
  $t_0(\mathbf 1)=N\p/2$, volumes) as the combinatorial error net;
  bit-match regression of the generalized $\mathbb{F}_5$ path.
\end{itemize}

% ---------------------------------------------------------------------
\section{The guaranteed lower bound at level $\p$}\label{sec:G1p}
\note{Source: CONGRUENCE.md \S7 (Theorem G1p).  Keep short: state the
deltas vs Paper I's Theorem G1 and prove only what is new.}
\begin{itemize}
\item Lemma D0 per cusp (exact trace reproduction; cross-sections tiled
  by mesh faces per copy) $\Rightarrow$ uninflated rank-two boundary
  block, $\sigma_h=\sqrt{N\p+1}\,\alpha_h^{\mathrm{ref}}$ (disjoint
  copies).
\item Multi-sliver absorption (per-copy reference constants); window
  data; the coefficient list; Theorem G1p statement + proof-by-delta.
\end{itemize}

% ---------------------------------------------------------------------
\section{Certificates and results}\label{sec:results}
\note{Sources: CONGRUENCE.md \S\S6, 8 status blocks;
theorem4postmortem.md (worktree).  One frozen-certificate table per
theorem; the $N\p=13$ subsection tells the stress-test story.}
\begin{itemize}
\item Frozen certificates (booktabs table): mesh, dofs, $\nu^*$,
  windows, $(\theta,\theta_2,\alpha,\theta_4,\tilde\rho)$, runtime,
  per-window $c_e$ vs $d_e$ ranges.
  [[$N5$: 8$\times$4$\times$3, 28{,}400 dofs, $\nu^*{=}1.02$, 8 windows,
  $(0.6,0.9,0.2,0.85,9.0)$.
  $N9$: 6$\times$3$\times$3, 26{,}532 dofs, $\nu^*{=}1.02$, 8 windows,
  $(0.5,0.9,0.15,0.8,1.5)$, $\sim$38 min.
  $N13$: 8$\times$3$\times$2, $\nu^*{=}1.001$, 16 windows,
  $(0.3,0.98,0.1,0.85,0.5)$, $\sim$93 min, 16/16.]]
\item Margin ladder and scaling discussion: $\mu$ vs volume
  (sublinear); the $\tilde\rho$ demand shrinking with $N\p$; the
  parameter-search methodology (scalar conditions + constant-direction
  heuristic + float pencil filter; floats choose, arb verifies; the
  scalar-slack/pencil anti-correlation through $\tilde\rho$).
\item The $N\p=13$ stress test: margin collapse to $\mu\approx 1.45$
  (mesh-converged); why the unscaled certificate fails while the matrix
  is SPD (stiffness-scale Rump shift vs mass-scale eigenvalue floor);
  resolution by Paper I's two refinements (error-free diagonal
  congruence + per-row radius absorption).  [[This subsection is the
  strongest evidence that the certificate, not the spectrum, was the
  binding constraint — and that the machinery scales.]]
\end{itemize}

% ---------------------------------------------------------------------
\section{Discussion and outlook}
\note{Source: CONGRUENCE.md \S8 ``beyond the ladder''.}
\begin{itemize}
\item Where the ladder ends as-is: the mode bound $N\p<4\pi^2Y^2$
  ($\approx 61$), certificate cost growth, and what each would take
  (taller $Y$/more windows; out-of-core verified Cholesky).
\item Composite levels ($>2$ cusps; general width bookkeeping);
  other Bianchi fields ($\ZZ[\omega]$, where the companion
  trace-formula engine of the parent project already has level-one
  data); torsion-free $\Gamma(\mathfrak n)$ covers.
\item Relation to the classical program: what certified gaps at small
  level do and do not say about the Selberg-type conjecture over
  $\mathbb{Q}(i)$.
\end{itemize}

% ---------------------------------------------------------------------
\section{Dependencies and reproducibility}
\note{Dependency delta vs Paper I: NONE (state and mean it — all new
ingredients are proved here or in Paper I).  Commands, runtimes,
software manifest; worktree/commit hashes for the certified runs.}

\appendix
\section{The cusp-$0$ audit}
\note{Import 0cuspchecks.md C1--C9 as a verification appendix ---
referee-friendly: every cusp-$0$ claim with its one-paragraph proof and
the numerical cross-checks that would catch each error class.}
\section{Gluing tables}
\note{The $\mathbb{P}^1(\ZZ[i]/\p)$ actions for all four generators at
$N\p=5,9,13$; cusp-class partitions.}

\end{document}
